\documentclass[a4paper,11pt]{article}

% ---------- Packages ----------
\usepackage[margin=2.5cm]{geometry}
\usepackage{amsmath, amssymb, amsthm, mathtools}
\usepackage{algorithm}
\usepackage{algpseudocode}
\usepackage{enumitem}
\usepackage{graphicx}
\usepackage{booktabs}
\usepackage{hyperref}
\usepackage{xcolor}

\usepackage{tikz}
\usetikzlibrary{automata,positioning,arrows.meta}

% ---------- Theorem Environments ----------
\newtheorem{theorem}{Theorem}
\newtheorem{lemma}{Lemma}
\newtheorem{definition}{Definition}
\newtheorem{proposition}{Proposition}
\newtheorem{corollary}{Corollary}

% ---------- Custom Commands ----------
\newcommand{\R}{\mathbb{R}}
\newcommand{\N}{\mathbb{N}}
\newcommand{\Z}{\mathbb{Z}}
\newcommand{\Alphabet}{\Sigma}
\newcommand{\eps}{\varepsilon}
\newcommand{\set}[1]{\left\{ #1 \right\}}
\newcommand{\abs}[1]{\left| #1 \right|}
\newcommand{\ceil}[1]{\left\lceil #1 \right\rceil}
\newcommand{\floor}[1]{\left\lfloor #1 \right\rfloor}

% ---------- Section Style ----------
\renewcommand{\thesubsection}{(\alph{subsection})}
\renewcommand{\thesection}{}

% ---------- Algorithm Style ----------
\algrenewcommand\algorithmicrequire{\textbf{Input:}}
\algrenewcommand\algorithmicensure{\textbf{Output:}}

% ---------- Title Info ----------
\title{\textbf{CSE2315 --- Assignment 5}}
\author{Vlad Paun \\ 6152937}
\date{\today}

% ---------- Document ----------
\begin{document}

	\maketitle

	\section{Exercise 1}
	Consider the alphabet $\Alphabet = \set{a,b,c}$ and the language
	\[
	L = \set{a^k b^k c^k \mid k \geq 1}.
	\]

	Give a high-level description for a three-tape enumerator that enumerates $L$. You may want to use the second and third tapes as counters. You do not have to give implementation details of the counters.


	\textbf{Answer: }Use Tape~1 as the printer and Tapes~2 and~3 as counters.
	Tape~2 stores the current value of $k$, initially $k=1$.

	In each iteration, the machine copies $k$ from Tape~2 to Tape~3 and uses
	Tape~3 to print $k$ copies of $a$, $b$, and $c$ on the printer tape,
	producing the string $a^k b^k c^k$. It then increments $k$ on Tape~2.

	Repeating this process prints $abc, aabbcc, aaabbbccc, \dots$, thus
	enumerating all strings in $L = \{a^k b^k c^k \mid k \ge 1\}$.



	\section{Exercise 2}
	Old exam question endterm 23--24:

	\subsection{Give a Turing Machine $M$ that is not a decider such that $L(M)$ is decidable.}

	Let $M$ be a Turing machine that accepts exactly the strings $\set{0^k1^k | k>0}$, and loops forever on every other input.

	\subsection{Explain why $M$ is not a decider.}

	Then $M$ is not a decider, as it does not halt on every input. It always either halts and accepts or loops forever.

	\subsection{Explain why $L(M)$ is decidable.}

	However, $L(M) = \set{0^k1^k | k>0}$ is decidable as shown in previous exercises.

	\section{Exercise 3}
	Old exam question resit 20--21: Let
	\[
	MIA_{TM} = \set{\langle M_1, M_2\rangle \mid L(M_1)\cap L(M_2)\neq \emptyset}.
	\]

	Show that this is a recognizable language.

	\textbf{Answer: }We construct a Turing machine $M$ that simulates $M_1$ and $M_2$ and runs every word in the alphabets language on both $M_1$ and $M_2$ (using dovetailing) and accepts if at any point the same word is accepted by both $M_1$ and $M_2$
	\vspace{1cm}

	\section{Exercise 4}
	In the game of chess, consider an intermediate board configuration, and suppose it is white's move. We know that white is either guaranteed to win or not. Let
	\begin{align*}
		L = \{w \mid w \text{ represents a board configuration, and white is guaranteed to win}&\\
		\text{if it is white's move and white plays optimally}&\}.
	\end{align*}


	Is $L$ recognizable, decidable, or neither? Explain.

	\textbf{Answer: }$L$ is decidable for it is finite. Chess has only a finite number of board configurations and $L$ represents a subset of those. We can simply create a Turing machine $M$ that has (withing it's configuration) all the words in $L$ hardcoded.

	\section{Exercise 5}
	Revision, old exam question, midterm 21--22: Given a DTM $M$ with $\Alphabet = \set{a,b}$ such that:
	\begin{itemize}
		\item $M$ has no transition to its $q_{\mathrm{reject}}$, i.e. the reject state is an unreachable state,
		\item $L(M) = \set{w \in \Alphabet^* \mid w \text{ starts with an } a}$.
	\end{itemize}

	Is $M$ a decider, not a decider, or can we not conclude that? Explain your answer.

	\textbf{Answer: }If an input word $w$ starts with a $b$, $M$ must never accept that word. But it also cannot go to it's reject state. So it must reject by looping. This means $M$ doesn't halt on every input so it is not a decider.

	\section{Exercise 6}
	Old exam question resit 20--21:

	In a log of a scheduler, it is noted down when a request for a task is given, and when an action is done to do that task. Requests are noted down with $r$, actions with $a$. A log is a final log if all tasks have been done. Thus, a final log can be described by the language of words where every $r$ has a corresponding $a$ that appears later in the string, formally given as
	\begin{align*}
		L = \{w \mid w &\text{ has an equal amount of $r$'s and $a$'s}\\
		&\text{and the $i$th $a$ appears later in the string than the $i$th $r$, for each } i \in (1, |w|/2)\}.
	\end{align*}


	Give an informal description of a Turing machine with the alphabet $\set{r,a}$ that decides this language. Use a maximum of 15 lines.

	As an example:
	\[
	\texttt{rarraa} \in L \qquad \text{and} \qquad \texttt{raaarr} \notin L
	\]

	\textbf{Answer: }This is just the valid parenthesis problem, but with $a$'s and $r$'s. Let Turing machine $M$:\\
	$M$ = "On input $w$:
	\begin{enumerate}
		\item Move to the end of the input.
		\item Write a special character $\$$ on the tape after the last input character and consider everything after the $\$$ as a stack.
		\item Go back to the first not crossed off character in the input and cross it off.
		\item If all input characters are crossed off, accept if the stack is empty, reject otherwise.
		\item If it's an $r$, push a special character $\#$ on the stack.
		\item If it's an $a$, pop from the stack. If the popped character is the $\$$, reject.
		\item Go to 3.
	\end{enumerate}

	\section{Exercise 7}
	Revision, old exam question midterm 19--20: Consider the following rules $R$ of a CFG
	\[
	G = (V,\Alphabet,R,S),
	\]
	with
	\[
	V = \set{S,A,B}:
	\]
	\begin{align*}
		S &\to ASA \mid Ba\\
		A &\to bB \mid S\\
		B &\to c \mid b \mid \eps
	\end{align*}

	\subsection{Give a valid set $\Alphabet$ such that $|\Alphabet| = 5$.}

	$\Alphabet=\set{a,b,c,d,e}$

	\subsection{Convert $G$ into Chomsky Normal Form (CNF). Use the procedure from Sipser, showing intermediate steps.}

	First, we add a new start variable.
	\begin{align*}
		S_0 &\to S\\
		S &\to ASA \mid Ba\\
		A &\to bB \mid S\\
		B &\to c \mid b \mid \eps
	\end{align*}
	Then, we remove the $\eps$-rules.
	\begin{align*}
		S_0 &\to S\\
		S &\to ASA \mid Ba \mid \textbf{a}\\
		A &\to bB \mid S \mid \textbf{b}\\
		B &\to c \mid b \mid \textcolor{gray}{\eps}
	\end{align*}
	Third, we remove all unit rules.
	\begin{align*}
		S_0 &\to \textcolor{gray}{S}\mid \textbf{ASA $\mid$ Ba $\mid$ a}\\
		S &\to ASA \mid Ba \mid a\\
		A &\to bB \mid \textcolor{gray}{S}\mid \textbf{ASA $\mid$ Ba $\mid$ a} \mid b\\
		B &\to c \mid b
	\end{align*}
	Convert the remaining rules into proper form.
	\begin{align*}
		S_0 &\to TA \mid BU \mid a\\
		S &\to TA \mid BU \mid a\\
		A &\to VB \mid TA \mid BU \mid a \mid b\\
		B &\to c \mid b\\
		T &\to AS\\
		U &\to a\\
		V &\to b
	\end{align*}


\end{document}